\documentclass[a4paper,14pt]{extarticle}
\usepackage{styles}

\begin{document}

\litSubsubsection{2.1.3 Prior GPU--GA Studies on GA Core Operations Optimization}

This subsection reviews relevant studies to identify patterns in GPU-accelerated GA computing. The review also examines the Python libraries and tools used.

\textbf{Study on Parallel Hybrid Genetic Algorithm and Machine Learning on GPU and Applications to Railway Scheduling.} Cheng and Gen~\cite{cheng2020parallel} parallelized a hybrid multi-objective GA applied to railway scheduling problems (ARL, BTS, and Japanese high-speed rail). The implementation offloaded differential evolution mutation, k-means clustering, and fitness evaluation to GPUs, which was necessary for maintaining practical computation times. The parallel implementation preserved solution quality while improving efficiency.

\textbf{GPU-Accelerated NSGA-II for Crystallization via Numba CUDA.} Rangavajhala et al.\ apply GPU-parallelized NSGA-II to crystallization process optimization, with the entire pipeline implemented using Numba CUDA kernels, achieving a $122\times$ speedup over the CPU baseline~\cite{rangavajhala2024accelerating}.

\textbf{GPU-Accelerated Stack-Based Genetic Programming: Sathia et al.} Candidate expressions are represented as prefix lists evaluated using fixed-length stacks in GPU memory. Benchmarked on an NVIDIA DGX-A100, the implementation achieves $119\times$ and $40\times$ speedups over gplearn on synthetic and large-scale datasets respectively~\cite{sathia2021accelerating}. The algorithm has been integrated into cuML, NVIDIA's scikit-learn-compatible GPU library~\cite{sathia2021accelerating}.

\textbf{EvoGP: GPU-Accelerated Tree-Based Genetic Programming.} Wu et al.\ present EvoGP, a GPU-accelerated framework for tree-based genetic programming that achieves up to $304\times$ speedup over existing GPU-based TGP implementations and $18\times$ over state-of-the-art CPU-based libraries~\cite{wu2025evogp}.

\textbf{GPU-Accelerated NSGA-III for Protein Encoding.} Kim and Kim report a $15{,}454\times$ speedup over Pymoo, enabling scaling from 100 solutions over 100 cycles to 400 solutions over 12,800 cycles~\cite{kim2024gpu}. The speedup magnitude partly reflects the choice of baseline (Pymoo is a Python framework). NSGA-III surpassed NSGA-II beyond 3,200 cycles, showing that GPU acceleration enables both faster computation and better solution quality through more extensive search~\cite{kim2024gpu}.

\textbf{Conclusions and implications.} The literature demonstrates that large speedups are achievable when GA operations are mapped to CUDA's thread hierarchy. Reported accelerations range from $18\times$--$122\times$ over CPU baselines~\cite{wu2025evogp,rangavajhala2024accelerating,sathia2021accelerating} to $304\times$ over prior GPU implementations~\cite{wu2025evogp} and $15{,}454\times$ over Python frameworks~\cite{kim2024gpu}, though speedup magnitudes depend on the choice of baseline. Fitness evaluation is the primary acceleration target, though parallel selection, crossover, and mutation can further reduce runtime. The dominant strategy involves CUDA C/C++ kernels with Python orchestration, whether through low-level bindings like PyCUDA or GPU-native libraries like cuML~\cite{sathia2021accelerating}. These findings inform this thesis's design, where fitness evaluation dominates GPU computation and genetic operators run as parallel kernels.

\end{document}
